\section{Discusión}
% Interpretacion propia de los datos
Como podemos ver en los horarios propuestos (Figuras~\ref{fig: lunes},~\ref{fig: martes},~\ref{fig: miercoles},~\ref{fig: jueves},~\ref{fig: viernes})
En todos los dias tenemos cruces de horarios debido a que a pesar de ser penalizados el AG siguio considerando estos los mejores resultados
segun nuestos indiviuos iniciales asi como nuestros Operadores Geneticos.

Con esto en mente es importante recordar que como se menciono en la seccion~\ref{ejecucion} usamos una población de 200 individuos iniciales
a lo largo de 500 generaciones debido a la gran combinacion de valores puede que esta poblacion sea insuficiente. Sin embargo se tomo esta decision
debido al poder computacional que tenemos. Por lo que para encontrar una mejor solucion se sugiere tener una poblacion mayor con mas generaciones,
teniendo en cuenta que esto tomara más tiempo de compilacion. En nuestro caso tomo alrededor de 90 min. O lo que es lo mismo 1 hora y media, por lo que
proponer al menos 500 individuos con 1000 generaciones tomaria aproximadamente el doble o triple de esto.

Sin embargo en nuestro caso practico podemos determinar que nuestro acercamiento es bastante aceptable ya que los casos donde existe cruce de salones
pueden ser resueltos con otros salones disponibles. Recordemos que el fin de los AG es encontar una solucion optima, tal vez no sea la ideal
pero si la mejor que tenemos con los recursos disponibles.